\documentclass[11pt]{article}

\usepackage[margin=1in]{geometry}
\usepackage{amsmath,amssymb}
\usepackage{booktabs}
\usepackage{longtable}
\usepackage{array}
\usepackage{xcolor}
\usepackage[hidelinks]{hyperref}

\setlength{\emergencystretch}{3em}

\definecolor{DoneGreen}{RGB}{20,115,70}
\definecolor{RunOrange}{RGB}{180,90,0}
\definecolor{PendRed}{RGB}{160,25,25}
\newcommand{\done}{\textcolor{DoneGreen}{\textbf{$\checkmark$ Done}}}
\newcommand{\running}{\textcolor{RunOrange}{\textbf{Running}}}
\newcommand{\pending}{\textcolor{PendRed}{\textbf{Pending}}}
\newcommand{\blocked}{\textcolor{PendRed}{\textbf{Blocked}}}
\newcommand{\bounded}{\textcolor{RunOrange}{\textbf{Bounded}}}
\newcommand{\authors}{\textcolor{RunOrange}{\textbf{Authors}}}

\title{\textbf{Experiment Progress and Next-Step Plan}}
\author{}
\date{13 September 2026 --- IEEE TIFS, fourteenth independent review cycle.
This file carries \emph{only} open, running and not-yet-started work;
completed items are removed once recorded in \texttt{docs/progress.md}.
Earlier editions are in \texttt{docs/archived/}.}

\begin{document}
\maketitle

\section*{Where This Cycle Stands}

The active review is the fourteenth independent review in
\texttt{docs/RevisionSuggestions.tex} (12 September 2026, IEEE TIFS, verdict
\emph{major revision}), findings R1--R14: five Major, seven Moderate, two
Minor. \textbf{All fourteen are closed, R3 in the review's strong form}: the
published mechanism it asks for was run end to end from its public release,
at matched delivered distortion, and does not separate from an isotropic
control of the same delivered energy. The cycle's substance is that two
of the review's Major findings changed what the paper claims rather than how
it says it, and both changed it against the authors: closing the attacker
panel broke the central null on one attacker of five, and the geolocator arm
produced a result the paper's own mechanistic account does not predict.

\paragraph{Compute state.} \textbf{Nothing is running; the vGPU~3090 was shut
down once the geolocator rows were home and verified.} The two arms that
needed a card are both finished. R1 was a re-embedding rather than a search
--- nine prescribed rules against three attackers that had never read them,
three seeds, same manifest, gallery, delivered distortion and energy gate ---
\textbf{36{,}000 rows over nine jobs}, all exiting zero. R4 was the geolocator
arm, \textbf{6{,}400 rows} after a defect described in Table~\ref{tab:tifs14-built}
forced the cheap arm to be rerun from scratch. The work done \emph{this}
sitting needed no GPU at all: the operator table R8 and R13 both ask for
recomputes from export trees already committed to the repository.

\paragraph{Document state.} \texttt{main.pdf} \textbf{13 pages} against the
13-page TIFS ceiling; abstract \textbf{245} words with \textbf{no math mode},
inside the SPS 150--250 window and its prohibition on equations;
\texttt{supplementary.pdf} \textbf{10 pages} against a 6-page guidance that is
a recommendation rather than a limit, with the Editor-in-Chief request and its
reasoning written into \texttt{docs/cover\_letter.txt};
\texttt{titlepage.pdf} 1 page; the extended evidence report, released with the
code and not submitted, 15 pages. All four build with \textbf{0 undefined
references, 0 undefined citations, 0 overfull boxes and 0 font-shape
warnings}. The abstract had drifted to \textbf{256} words when the geolocator
sentence went in --- over a hard venue cap, in a document nobody rebuilds the
word count for --- and six generated tables were running past the column by 9
to 44\,pt. Both classes are now checked by
\texttt{src/scripts/check\_build\_health.py} rather than by eye.

\paragraph{Verification state.} \textbf{902 claims verify against raw exports;
0 mismatched, 0 unverifiable.} The three added this cycle are the operator
arm's, re-anchored from the extended report to the supplement, plus the four
released-amplitude figures R7's own fix introduced --- which had replaced one
wrong number with four unregistered ones, so the sentence written to reconcile
a cross-document contradiction was itself unchecked for a cycle. All four are
correct: they are means over both attackers' rows, which is why they differ in
the first decimal from the per-attacker figures the manuscript quotes beside
them. \textbf{24 registered values are no longer printed anywhere.} That count
is not a defect and not noise to be silenced: every one verifies, and each was
compressed out of prose in an earlier cycle while its table twin still prints
it. It is kept visible because the alternative --- deleting a true measurement
from the registry when the page budget squeezes it out of a sentence --- is
the wrong direction.

\paragraph{A checker that could not tell a drop from a table.} The claim
auditor asks whether a claim's printed string still appears in the document
meant to print it, and it read that document literally, without following
\verb|\input|. Every value reaching the page through a generated table
therefore answered ``no''. The genuine drops were indistinguishable from the
false ones, which is the single thing this check exists to prevent. It now
splices one level of \verb|\input| before looking.

\paragraph{A pointer to a table that did not exist.} \S III-E of the
manuscript said its three operator results were ``tabulated per operator in
the Supplementary Material''. The supplement carried the operator
\emph{definitions} and no results table: the sentence had been true of the
extended report, which is not submitted, and stayed on the page when the
tables moved. A referee checking the claim that the null survives every
operator would have found a sentence and nowhere to check it. The table now
exists, built from the committed operator exports, and it
reproduces every cell the claim registry already held --- which is how a
missing table went eight cycles without being noticed as missing.

\section*{Table 1: Open Work Against the Fourteenth Review}

\begin{longtable}{@{}p{0.05\linewidth}p{0.19\linewidth}p{0.11\linewidth}p{0.57\linewidth}@{}}
\caption{All fourteen findings. R1 and R4 changed what the paper claims; R3 is
closed in the strong form, by running the published mechanism end to end from
its public release.}
\label{tab:tifs14-open}\\
\toprule
\textbf{ID} & \textbf{Item} & \textbf{Status} & \textbf{Outcome} \\
\midrule
\endfirsthead
\toprule
\textbf{ID} & \textbf{Item} & \textbf{Status} & \textbf{Outcome} \\
\midrule
\endhead

R1 & Two axes read by different attacker panels & \done & The finding the
review called its strongest counter-argument, and it was right. The two
attackers a solved map separates on --- Patch-NetVLAD and CLIP --- were
exactly the two no prescribed rule had ever faced. Extending the same nine
rules to Patch-NetVLAD, ViT-B/16 and CLIP ViT-L/14 broke the null on one of
them: \textbf{four of nine separate on Patch-NetVLAD} (edge magnitude
$-0.0550$, score gradient $-0.0492$, margin gradient $-0.0308$, and the other
way, anti-margin gradient $+0.0208$), none on the ViT or CLIP. Two facts
outrank the count and are in the paper: the largest winner is \emph{edge
magnitude}, the same rule that beats uniform once the budget passes the
operating point, so edge alignment is the one prescribed strategy that appears
wherever allocation works and it is not a sensitivity map; and the
\emph{margin gradient}, the rule this paper derives rather than borrows,
separates while its inverse is significantly worse --- the account's first
confirming case. Why Patch-NetVLAD we cannot say, and the paper says so.
The five-attacker placement panel is now a supplementary table; Fig.~2 redrawn at five attackers; abstract, contributions,
\S III-C and conclusion rewritten. \\
\addlinespace

R2 & Abstract stated the null above what the intervals certify & \done &
Rewritten as a bounded non-detection: three of nine non-degenerate cells
certified at the declared margin, the rest undecided, and the certifying
margin for all nine printed as $\pm0.027$. TOST columns added to both
placement tables; the number of query clusters certifying $\pm0.01$ would take
--- $1{,}261$ weak, $917$ strong, two to three times this corpus --- is
printed, which turns a limitation into a design specification. \\
\addlinespace

R3 & No published mechanism run end to end & \done & Closed in the review's
strong form. GeoShield's public release was run end to end under this paper's
protocol --- same attacker, same gallery, same delivered distortion --- over
400 query clusters at one seed, against two controls: no perturbation, and an
isotropic perturbation of identical delivered energy. Against that
energy-matched control it \textbf{does not separate}: $-0.0100$ Top-1,
place-clustered interval $[-0.041,+0.020]$, $p=0.52$ over 38 discordant pairs.
At one seed that interval is wider than the $\pm0.01$ margin the other arms
certify against, so this is a \emph{bounded null}, not a certified
equivalence, and it is reported as such. Clearing the release took five
obstacles, four found by reading and one only by running; they are recorded in
the supplement as a reproducibility observation about that artefact. Because
the release stubs its vision--language call to a constant caption, the arm
measures \textbf{the released artefact, not the method as described}, and the
manuscript, the provenance table and this row all say so in the same sentence
so neither half travels alone. Not done: running it with a working
vision--language component, which is N1. \\
\addlinespace

R4 & Geolocators excluded from the threat model & \done & The arm ran, and it
produced the one result this paper could not have predicted. \emph{Direction
transfers; allocation does not.} On the 218 queries the clean attacker
localises within 25\,km: isotropic $0.713$, direction $0.546$ ($-0.167$
$[-0.224,-0.107]$), hardened $0.514$ ($-0.199$ $[-0.260,-0.137]$), Holm
corrected; edge $+0.040$ and saliency $+0.023$, neither separating. The paper
explains the axis asymmetry through a correct-versus-hardest-negative
\emph{ranking} margin --- and a model emitting a coordinate has no gallery, no
hardest negative and no rank to flip, so the manuscript had said in as many
words that neither half of the argument transfers. The explanation does not
transfer and the effect does. That is now stated as a fact about the finding
and a limit on the explanation, not smoothed over. The geolocator table is in the supplement; the
abstract's and conclusion's ``untested'' sentences are gone. \\
\addlinespace

R5 & Uncorrected exploratory family, no equivalence test, unjustified test &
\done & All four parts. The placement family is now confirmatory and Holm
corrected within itself --- correction changes no verdict, which is worth more
than if it had. TOST at the declared margin runs on every placement contrast.
A place-clustered sign-flip permutation test, which assumes nothing about the
shape of a three-seed-averaged binary outcome, agrees with the signed-rank
column in every cell. And ``1,200 paired observations'' is corrected to the
400 differences the declared unit actually produces, with the seed averaging
stated in the same sentence so the two sections cannot be read apart. \\
\addlinespace

R6--R7 & A miscount, a mis-transcribed interval, two irreconcilable
amplitudes & \done & Nine contrasts are reported as nine everywhere; the
quoted interval is edge magnitude's $[-0.014,+0.025]$ with its unit named,
where the old one took a lower bound from one column and an upper from
another. The $49.30$-versus-$76$ contradiction was resolved by recomputing:
the $76$ matched no row, and the four figures that replaced it are named by
condition --- $12.2$ and $84.0$ at the operating point for uniform and edge
placement, $48.7$ and $241.4$ at the large budget. Those four were unregistered
until this sitting. \\
\addlinespace

R8 & Evidence deferred outside the submission & \done & The last dangling
pointer is closed, and it was the worst kind: a claim of a table that did not
exist (see above). The operator arm is now a supplementary table. Thirteen ``extended
report'' pointers are down to eight, each naming what it defers; four tables
remain outside and are listed once at the head of the supplement. \textbf{No
claim in the abstract or the conclusion now rests on a table absent from the
submitted package}, which is the review's acceptance criterion verbatim. \\
\addlinespace

R9 & IEEE SPS and TIFS format compliance & \done & Abstract 245 words, no math
mode --- it was 256 and over a hard cap when this cycle started, having drifted
there one sentence at a time. Main 13 of 13. Supplement 10 pages with the
request and its reasoning in the cover letter, which is what SPS guidance
permits rather than a violation to be hidden. Acknowledgment, IEEE
generative-AI disclosure and a titled ethics and data-availability section all
present. Biographies are written and commented out until acceptance, since
IEEE counts them against the 13-page limit and the space is spent on evidence. \\
\addlinespace

R10 & TIFS deep-learning reproducibility checklist & \done & All five gaps, and
one of them was a finding rather than an omission: the released checkpoint was
trained with \emph{no mask source}, and the training path substitutes a
constant zero field when none is given, so that run optimised against a target
carrying no spatial information at all. That is the cause of the degeneracy the
paper leads with, and it is recoverable from the checkpoint --- weights moving
$6.8\times10^{-3}$ across five epochs. Weight decay, constant learning rate,
absence of normalisation layers, the mask-supervised corpus and its class
balance are all now stated, with a paragraph mapping the checklist item by
item. \\
\addlinespace

R11--R12 & Missing literature; cross-corpus admissibility & \done & The
protective-perturbation-fails line and the adversarial-evaluation methodology
line are both cited, each with what this paper adds rather than a bare
citation. Admissibility was recomputed on the 400 MSLS query frames the attack
actually sees: the \emph{ordering} is the VOC ordering at every one of twelve
cells, so which axis costs more does not depend on the corpus, while the
\emph{level} is not comparable. The VOC-calibrated tolerance is therefore not
carried across, and Table~IV's caption says where it was calibrated and what
that assumes. \\
\addlinespace

R13 & Presentation for a TIFS readership & \done & Conventional signposts added
as an opening paragraph of \S III-A naming the setup, controls, baselines and
ablations --- which also carries R1's requirement that the attacker panel be
stated once per axis instead of reconstructed from three sections. The settled
position on allocation is now stated at the head of \S III-C \emph{before} the
evidence that produced it, with the discovery order marked as such. The
operator arm is tabulated (R8); KITTI-360 already was. The mechanism diagram is
restored but sits in the supplement: at 13 of 13 pages the main body keeps the
protocol schematic instead, which draws this paper's own contribution rather
than its testbed. \\
\addlinespace

R14 & Residual items & \done & All six. The gallery-ceiling extrapolation now
states its specific assumption and when it fails, rather than calling web
scale ``an extrapolation''. The negative inferences drawn from five attackers
are restated as what they are --- no pattern apparent across five, and five too
few to support a claim either way. The $\pm0.01$ margin is described as
declared in advance in the released repository rather than in a third-party
registry, and this cycle that honesty was propagated into the figure caption
and both table captions, which still read ``fixed before testing''. The
bibliography is 39 entries with 0 uncited and 0 cited-but-missing. The question
title is kept deliberately. \\

\bottomrule
\end{longtable}

\section*{Table 2: How This Cycle's Arms Were Built, and What Went Wrong}

\begin{longtable}{@{}p{0.16\linewidth}p{0.78\linewidth}@{}}
\caption{Construction decisions and defects, recorded because none is
recoverable from the result. Two of the four are self-inflicted bugs.}
\label{tab:tifs14-built}\\
\toprule
\textbf{Aspect} & \textbf{Detail} \\
\midrule
\endfirsthead
\toprule
\textbf{Aspect} & \textbf{Detail} \\
\midrule
\endhead

R1, the gate that
makes it a re-embedding &
Extending nine rules to three new attackers is only a re-embedding if it is
the \emph{same} released frames being read by a different model, and that has
to be shown rather than asserted. It is: the objective trajectory is
$2.8412\to2.4645$ on every arm, bit for bit, and the top-decile concentration
is $0.547$ on all four sharing a map. Delivered MSE is $15.6800$ on every
perturbed condition and the energy gate reads $1.000000$ throughout. Without
those checks the panel would be a new search wearing the old one's name. \\
\addlinespace

R4, why the
contrast is read on a subset &
The geolocator is scored on the 218 of 400 queries the \emph{clean} attacker
already localises within 25\,km, and the reason is concrete rather than
statistical. In the smoke run GeoCLIP placed a Boston frame in Paraguay, 7{,}654\,km
out --- and edge placement ``improved'' it to 1.09\,km. On a query the attacker
fails unaided, a defense collects credit it did not earn and damage it did not
cause. Both readings are reported; the subset is the one the claim rests on.
The 1\,km threshold moves for nothing, but at $6.0\%$ clean accuracy that is a
floor, not a finding, and is reported as such. \\
\addlinespace

R4, a race that
exposed a real bug &
Hand-packing three cheap jobs into idle workers hit the launcher's
\texttt{pgrep} race: two processes wrote one file, 8{,}541 rows where 6{,}400
were expected. Of 2{,}141 duplicate keys, 2{,}140 agreed bit for bit --- and
the single conflict was the one that mattered. The runner seeded its generator
per \texttt{(query, seed)} and advanced it across conditions, while the resume
path skipped completed conditions with a bare \texttt{continue} \emph{without}
advancing the generator. A resumed run and a fresh run therefore drew
different noise fields, which violates this repository's own rule that a
long-running script must resume to the same state it would have reached
uninterrupted. Reseeded per \texttt{(query, condition, seed)} and the whole
cheap arm rerun into a fresh directory under a single writer: 6{,}400 rows,
6{,}400 unique keys, delivered MSE exactly $15.68$ per condition. \emph{The
duplicate was noise; what it uncovered was not.} \\
\addlinespace

R8/R13, the operator
table &
No GPU and no new run: the operator study's per-query rows have been committed
since the ninth cycle, so the missing table was only ever missing. It is built
at the protocol's own unit --- three seeds averaged within a query, 400 paired
differences, intervals bootstrapped over the 277 place clusters --- rather
than at the 1{,}200-row unit the archived analysis CSVs use. Every cell
reproduces the values already in the claim registry, which is the check that
the new generator reads the data the same way the old analysis did. \\
\addlinespace

The KITTI-360 pairing,
carried forward &
Still the construction behind the second-dataset claim. A place is a ground
cell the drive enters twice, the later visit giving the query and the earlier
the positives, every pair at least 600 frames apart, with adjacent revisited
cells merged first --- two cells metres apart are one street corner and, left
separate, generate each other's false negatives. Honest limit, unchanged: 16
places behind 227 queries, so the clustered interval is $1.93\times$ the query
one and covers zero for direction. The allocation null is unaffected. \\

\bottomrule
\end{longtable}

\section*{Table 3: Author-Side Items}

\begin{longtable}{@{}p{0.05\linewidth}p{0.26\linewidth}p{0.10\linewidth}p{0.51\linewidth}@{}}
\caption{Not manuscript edits, experiments or analyses. Carried forward; S6 is
new this cycle.}
\label{tab:tifs14-authors}\\
\toprule
\textbf{ID} & \textbf{Item} & \textbf{Status} & \textbf{What is needed} \\
\midrule
\endfirsthead
\toprule
\textbf{ID} & \textbf{Item} & \textbf{Status} & \textbf{What is needed} \\
\midrule
\endhead

S1 & ORCIDs, EDICS, submission metadata & \authors & Verifiable only in the
submission system, not in the files. All three authors need ORCIDs registered
and the EDICS categories chosen; the metadata must match the manuscript's
title and author list. \\
\addlinespace

S2 & Prior-submission disclosure & \authors & The repository holds a response
letter to an earlier venue. Whether disclosure is required depends on facts
outside these files; if it is, the statement must be made truthfully at
submission. \\
\addlinespace

S3 & Responsible-use and data-terms wording & \done\ (read) & Now a titled
section of the manuscript rather than the closing paragraph of a theory
section (R9). Wording confirmed read; re-wrapped since for the page budget. \\
\addlinespace

S4 & DOIs in the bibliography but not in the PDF & \done\ (decided) &
\texttt{ref.bib} carries verified DOIs; the mandated \texttt{IEEEtran.bst}
v1.14 has no \texttt{doi} support, so none is typeset. Decision: submit as is. \\
\addlinespace

S5 & Response letter & \done\ (not required) & Confirmed not being written. \\
\addlinespace

S6 & Supplement over the six-page guidance & \authors & \textbf{Needs the
author's sign-off before submission, not more work.} The supplement is 10
pages against a 6-page recommendation, and the request plus its reasoning are
written into \texttt{docs/cover\_letter.txt}: a negative result against an
assumption the field designs around is held by TIFS's own deep-learning
guidance to a higher reproducibility standard, and a referee testing one needs
the cells it was pooled from. The letter also states what happens if the
Editor-in-Chief declines --- the four tables go back to the extended report and
the manuscript says so. Confirm you are content to ask. \\

\bottomrule
\end{longtable}

\section*{Table 4: Not Scheduled}

\begin{longtable}{@{}p{0.05\linewidth}p{0.26\linewidth}p{0.10\linewidth}p{0.51\linewidth}@{}}
\caption{Not required for this submission. Listed so a later cycle does not
rediscover it. N1 is the largest remaining scientific gap.}
\label{tab:tifs14-optional}\\
\toprule
\textbf{ID} & \textbf{Item} & \textbf{Status} & \textbf{Why it is not blocking} \\
\midrule
\endfirsthead
\toprule
\textbf{ID} & \textbf{Item} & \textbf{Status} & \textbf{Why it is not blocking} \\
\midrule
\endhead

N1 & A published mechanism run end to end \emph{with its vision--language
component supplied} & \pending & Narrowed, not blocked. The public release now
runs end to end here (R3); what remains is supplying the component it stubs,
which would measure the method as described rather than the artefact as
shipped. The runner already accepts a caption source and caches per target, so
this is a budget decision rather than an engineering one --- scoped at roughly
one further GPU day and deliberately not bought, because the released-artefact
arm answers the review's question and the substituted-component arm answers a
different one. Reopen if that distinction becomes load-bearing for a
reviewer. \\
\addlinespace

N2 & A stronger geolocator than GeoCLIP & \done & Run, with a substitute,
because the model the review names has no public weights: PIGEON is not
released, so the strongest geolocator obtainable is the strongest
\emph{public} one. We prompt a vision--language model for a coordinate --- a
different architecture, different training data, no gallery --- which also
tests whether the direction result is a property of the CLIP family that every
other attacker here shares. It clears the capability gate on all 400 clean
frames: $65.0\%$ within 25\,km at a median of $6.1$\,km against GeoCLIP's
$54.5\%$ and $18.5$\,km, and nearly twice as accurate at street level, with no
refusals. \textbf{The result is two-sided and both halves are in the paper.}
The ordering survives: on the 260 queries it localises, the direction release
costs it $0.100$ more than matched-energy isotropic ($[-0.156,-0.046]$,
$p<0.001$), so the finding is not a CLIP artefact. The magnitude does not:
GeoCLIP retains $55.5\%$ of its localisable queries under the release and this
attacker retains $70.8\%$, with median error on that subset moving $3.4$ to
$5.3$\,km against GeoCLIP's $5.4$ to $16.7$. Against the best geolocator we
could obtain the release degrades city-level localisation without preventing
it, and \S Limitations now says so. Reported as a bound on what is currently
obtainable, not a ceiling: a stronger attacker may exist behind an interface
we cannot audit. \\
\addlinespace

N3 & Adaptive attacker at a larger budget & \done & Run, and it changed a
conclusion. The published arm fine-tunes the last residual block against a
fixed index; this runs the fully unfrozen encoder (every residual stage plus
the stem) beside it, both budgets in one batch under one protocol over four
seeds, each adapted model scored on the held-out split its own checkpoint was
trained against and differenced against an unadapted baseline on that same
split. \textbf{Against the direction release the larger budget helps, but only
when the attacker also re-embeds the gallery}: 12 of 12 cells positive across
three training exposures and four seeds, mean $+0.0625$ Top-1, sign-test
$p=0.00024$. Held at the pretrained index the same budget does nothing --- 4
of 12 positive, mean $-0.0050$, $p=0.93$. In level terms the unadapted
attacker reaches $0.05$--$0.08$ and the unfrozen re-indexing one $0.07$--$0.17$,
recovering roughly half of what the perturbation removed. The manuscript's
adaptive-attacker section therefore needs qualifying rather than extending:
its ``adaptation buys the attacker little'' holds for the budget it tested and
not for the one the review asked about. Note the split that matters is not
capacity but \emph{re-indexing} --- extra parameters are worth nothing to an
attacker forced to keep the original index. Two seeds could not have
established this: several cells flipped sign between seeds with non-overlapping
intervals, so the two-seed run would have supported a false positive or a false
null depending on which cells were read. \\
\addlinespace

N4 & Six segmenters re-scored on MSLS query frames & \pending & What it would
take to carry a calibrated admissibility tolerance onto the frames the attack
sees, rather than only the ordering (R12). The manuscript says exactly this is
what is missing and declines to carry the VOC tolerance across. \\
\addlinespace

N5 & Registry entries with no printed counterpart & \pending & 24 verified
values are no longer quoted in either document, having been compressed out of
prose while their table twins still print them. Nothing is wrong with any of
them; the open question is only whether the auditor should name this class
separately from a genuine drop. Kept visible rather than deleted. \\

\bottomrule
\end{longtable}

\section*{Closed in Earlier Cycles}

The tenth review's R1--R7 and the thirteenth's R1--R12 are closed, with the
evidence in \texttt{docs/progress.md}. The two worth carrying forward as
cautionary notes are both about checks that passed while the thing they
checked was wrong: a released export that had been appended to twice, whose
summary was computed by the same script that wrote the rows and so could not
detect the duplication; and three cross-document pointers that went stale
because \LaTeX{} cannot see a reference into a separately compiled document,
so the build reported zero undefined references while the manuscript pointed
at the wrong table. This cycle adds a third of the same kind: a claim registry
that verified 899 numbers while the sentence introducing four of them had
never been registered at all.

\end{document}
